% =============================================================================
%  CLAS12 Technical Note --- body
%
%  Section order follows the CLAS12 note authoring guide (CLAUDE.md, section 2).
%  Replace the bracketed [ ... ] guidance with the note content, and delete the
%  sections that do not apply.  Standard simulation constants are pre-filled from
%  the authoring guide (section 3); state any deviation explicitly.
% =============================================================================

\section*{Overview}
\addcontentsline{toc}{section}{Overview}
[One paragraph stating the single question the note answers: the background rate,
occupancy, PMT current, dose, or infrastructure change being studied, and why it
matters. State the answer in the Conclusions.]

\section{Experiment Configuration}
[Measurement notes only. Beam current, magnet states (solenoid, torus), which
detectors were on or off, and the run date. Reference a figure of the beam on
target when available. Omit this section for pure simulation studies.]

\section{Simulated Detector, Beam, and Target Configuration}
GEMC starting configuration: [FTOn / FTOff]. The target, beamline, shielding,
torus, and vacuum-beamline geometry are imported into Geant4 directly from the
engineering CAD models~\cite{note:cad,note:moller,note:beamline,note:target}.

A total of [124{,}000] electrons are propagated through the target within a
[248.5]~ns time window, corresponding to the full nominal CLAS12 operating
luminosity of $1\times10^{35}$~cm$^{-2}$\,s$^{-1}$ on a 5~cm liquid-hydrogen
(LH$_2$) target. [State the pair actually used and any survey offset.]

The beam energy is 10.6~GeV (results at 6.4~GeV agree to within ${\sim}10\%$).
The solenoid and torus are at full current; both torus polarities are reported as
in-bending and out-bending. Most rates show negligible (${<}{\sim}10\%$)
polarity dependence.

Software: GEMC [version] with Geant4 [version]. Physics lists: standard Geant4
electromagnetic; hadronic FTFP\_BERT; optical processes [included / not included].
[State the run-number scheme: legacy RUNNO=11 with \_mc CCDB variations, or a
run-dependent SQLITE geometry database with the \texttt{default} variation.]

\section{Results}
Report one subsection per detector or per studied quantity. For each: the
rate, occupancy, or current; the applied energy-deposition threshold; the
field-polarity dependence; and the comparison to data when available.

\subsection{[Detector] rates}
[Rate / occupancy / current, the applied threshold, the field-polarity
dependence, and the comparison to data if available. Report rates in kHz (MHz for
the largest), currents in $\mu$A, energy in MeV, and doses in rad/h. State the
threshold in every figure caption.]

An example rate distribution is shown in Fig.~\ref{fig:example}, and representative
values are collected in Table~\ref{tab:example}.

\begin{figure}[htbp]
  \centering
  \framebox[0.62\textwidth][c]{\rule[-1.4cm]{0pt}{3cm}\itshape figure placeholder}
  \caption{What is plotted, the configuration (FTOn/FTOff, in-bending/out-bending),
  the applied energy-deposition threshold, and a note that it is at the full nominal
  CLAS12 operating luminosity when applicable.}
  \label{fig:example}
\end{figure}

\begin{table}[htbp]
  \centering
  \begin{tabular}{lrr}
    \toprule
    Detector & Rate at 0 threshold (MHz) & Rate at 1~MeV (kHz) \\
    \midrule
    CTOF & 5.0 & 150 \\
    ECAL & 2.5 & 140 \\
    FTOF (Panel-2) & 6.0 & 700 \\
    \bottomrule
  \end{tabular}
  \caption{Example rate table. State the threshold for every quoted rate and the
  field polarity; remind the reader how to obtain the total (e.g. CTOF: multiply
  the per-counter rate by 48 counters).}
  \label{tab:example}
\end{table}

\section{Conclusions}
[The answer to the note's question, as numbered findings when there are several
distinct conclusions. Acknowledge any simulation-versus-data discrepancy rather
than hiding it.]
